\documentclass[12pt]{article}
\usepackage[T1]{fontenc}
\usepackage[utf8]{inputenc}
\usepackage{lmodern}
\usepackage{textcomp}
\usepackage{enumitem}
\usepackage{geometry}
\geometry{margin=1in}
\begin{document}
\begin{center}
Answer Key - Biology - Graduate Level Assessment 9 \
\small (Answers are ordered exactly as in the question paper.)
\end{center}

\section*{Multiple Choice}
\begin{enumerate}[label=\arabic*.]
\item \textbf{Answer:} A
\\ \textit{Options:} A) Smaller cells are more resilient to external damage.; B) Smaller cells have a larger number of organelles.; C) Smaller cells require less energy to function.; D) Smaller cells can multiply faster.; E) Smaller cells fit together more tightly.
\\ \textit{Note:} Smaller cells have a higher surface area-to-volume ratio, which enhances their ability to efficiently exchange materials and respond to environmental changes, making them more resilient to external damage compared to larger cells.
\item \textbf{Answer:} E
\\ \textit{Options:} A) disruptive selection; B) sexual selection; C) bottleneck effect; D) non-random mating; E) stabilizing selection
\\ \textit{Note:} Stabilizing selection describes the process by which intermediate traits, such as an average birth weight of 3 kg, are favored in a population because they provide a survival advantage, resulting in reduced variability and promoting the maintenance of this trait over time.
\item \textbf{Answer:} A
\\ \textit{Options:} A) Through vertical gene transfer from parent to offspring; B) Through random assortment of genes during meiosis; C) Through peptide bond formation errors during translation; D) Through genetic mutation only; E) Through transcriptional errors in RNA processing
\\ \textit{Note:} The correct answer is based on the principle that vertical gene transfer allows for the inheritance of genetic variations from parent organisms, which can lead to the production of diverse proteins through processes such as mutation and recombination during reproduction.
\item \textbf{Answer:} E
\\ \textit{Options:} A) By using a controlled group of dogs not exposed to the sounds; B) By using sensory deprivation to heighten hearing sensitivity; C) By observing the dog's natural behavior in different sound environments; D) By using negative reinforcement only; E) By using operant conditioning
\\ \textit{Note:} Operant conditioning allows researchers to reinforce a dog's response to specific sound frequencies, enabling them to measure the animal's sensitivity to those frequencies through behavioral changes in response to auditory stimuli.
\item \textbf{Answer:} A
\\ \textit{Options:} A) -.27, 1.45, 1.73; B) -.27, -1.45, -1.73; C) -2.18, 0.15, -1.82; D) 0, 1.5, -2; E) .27, 1.45, 1.73
\\ \textit{Note:} The correct answer is derived by using the Z-score formula, Z = (X - μ) / σ, where X is the raw score, μ is the mean (118), and σ is the standard deviation (11), resulting in Z-scores of approximately -0.27 for 115, 1.45 for 134, and 1.73 for 99.
\end{enumerate}

\section*{True / False}
\begin{enumerate}[label=\arabic*.]
\setcounter{enumi}{5}
\item \textbf{Answer:} True
\\ \textit{Options:} A) True; B) False
\\ \textit{Note:} Chemical shift subtraction can contribute to routine dynamic contrast subtraction in morphologic analysis particularly for the evaluation of margins of benign lesions in fatty breasts. It can also help in morphologic analysis of masses in dense breast.
\item \textbf{Answer:} True
\\ \textit{Options:} A) True; B) False
\\ \textit{Note:} It can be concluded that CLBP patients in general have worse motor task performance compared to healthy subjects and that provoking pain-related cognitions further worsened performance.
\item \textbf{Answer:} True
\\ \textit{Options:} A) True; B) False
\\ \textit{Note:} For a given NIHSS score, the median volume of right hemisphere strokes is consistently larger than the median volume of left hemisphere strokes. The clinical implications of our finding need further exploration.
\item \textbf{Answer:} False
\\ \textit{Options:} A) True; B) False
\\ \textit{Note:} Consistent with prior studies, we found an inverse relationship between obesity and serum prostate specific antigen. However, the magnitude of the difference was small. Thus, adjusting prostate specific antigen for body mass index does not appear warranted.
\item \textbf{Answer:} False
\\ \textit{Options:} A) True; B) False
\\ \textit{Note:} Our results suggest that colonoscopy does not affect the management of patients with acute diverticulitis nor alter the outcome. The current practice of a routine colonoscopy after acute diverticulitis, diagnosed by typical clinical symptoms and CT needs to be reevaluated.
\end{enumerate}

\section*{Long Form Response}
\begin{enumerate}[label=\arabic*.]
\setcounter{enumi}{10}
\item \textbf{Answer:} Interdisciplinary collaboration is crucial for a comprehensive understanding of the biosphere, particularly in the context of mapping climatic and vegetation patterns. Climatologists, ecologists, and geologists each bring unique perspectives and expertise that enhance our understanding of these complex interactions. For instance, climatologists often find it more beneficial to collaborate with geologists when investigating how geological formations influence climatic conditions and, consequently, vegetation patterns. This partnership allows for the integration of climatic data with geological insights, leading to more accurate models of ecosystem dynamics and better predictions of how climate change might impact biodiversity. Such collaborative efforts ultimately foster a holistic view of the biosphere, essential for effective environmental management and conservation strategies.
\\ \textit{Note:} correct because it emphasizes the necessity of interdisciplinary collaboration among climatologists, ecologists, and geologists to effectively analyze and interpret the intricate relationships between geological formations, climate, and vegetation patterns in the biosphere.
\item \textbf{Answer:} Chemical transmission at synapses is characterized by its unidirectional nature, primarily ensured by the localization of neurotransmitters and the structural organization of synapses. Neurotransmitters are synthesized and stored in synaptic vesicles at the axonal terminals, which are specifically designed for release into the synaptic cleft upon neuronal activation. In contrast, dendrites lack synaptic vesicles, preventing the uptake of neurotransmitters and ensuring that signals propagate in a single direction-from the presynaptic neuron to the postsynaptic neuron. This structural asymmetry, along with mechanisms such as receptor specificity and reuptake processes, reinforces the unidirectional flow of information in neural circuits. Thus, the combination of neurotransmitter localization at axonal ends and the specialized architecture of synapses plays a crucial role in maintaining the integrity of synaptic transmission.
\\ \textit{Note:} correct because it emphasizes that the unidirectional transmission of signals at synapses is facilitated by the specific localization of neurotransmitters in axonal terminals and the absence of synaptic vesicles in dendrites, which collectively prevent retrograde signaling.
\end{enumerate}

\vfill
\noindent\textit{End of Answer Key}
\end{document}